% AA-49 | Kwadratowy przypadek algorytmu Kovacica
% Klasyczne listy lokalne Ec i Einfty są cytowane z Kovacica (1986).
% Samodzielne dowody: symetryczny kwadrat, nieliniowy warunek śladu,
% odzyskanie obu pierwiastków Riccatiego i liniowość równania pomocniczego.
\chapter{Kwadratowy przypadek algorytmu Kovacica}
\chaptermark{Kwadratowy przypadek Kovacica}
\label{aa:ch:kovacic-quadratic}

W rozdziale~\ref{aa:ch:kovacic-rational-potential}
otrzymaliśmy kompletny algorytm znajdujący
wymierne rozwiązania równania Riccatiego
$r'+r^2=q$ dla każdego $q\in\mathbb C(x)$.
Nie rozstrzyga to istnienia rozwiązań
algebraicznych stopnia drugiego nad
$\mathbb C(x)$, co pokazuje już
Przykład~\ref{aa:prz:kovacic-quadratic-example}.
W niniejszym rozdziale badamy
\emph{drugą gałąź algorytmu Kovacica}:
rozwiązania Riccatiego w kwadratowych
rozszerzeniach $\mathbb C(x)$.
Przygotowujemy równanie pomocnicze
trzeciego rzędu i wykonujemy wszystkie
rachunki odzyskujące rozwiązania.
Skończoność katalogu możliwych
lokalnych danych wejściowych
opiera się natomiast na
\emph{klasycznym twierdzeniu Kovacica},
które jawnie cytujemy; nie przedstawiamy
jego ogólnego dowodu w tym rozdziale.

\section{Ślad i wyróżnik dwóch rozwiązań Riccatiego}

\begin{definicja}[Przypadek kwadratowy Riccatiego]
\label{aa:def:kov2-quadratic}
Niech $K=\mathbb C(x)$, $D=\mathrm d/\mathrm dx$,
a $q\in K$. Rozwiązanie $r$
równania $Dr+r^2=q$ w rozszerzeniu
algebraicznym $L/K$ nazywamy
\emph{ściśle kwadratowym}, jeśli
$[K(r):K]=2$. W takim przypadku
oznaczamy przez $\bar r$ drugi
pierwiastek jego wielomianu minimalnego.
Wówczas $\bar r$ również spełnia
równanie Riccatiego: jedyne
przedłużenie $D$ do
separowalnego rozszerzenia
algebraicznego komutuje z
automorfizmem sprzężenia.
\end{definicja}

\begin{lemat}[Ślad, iloczyn i wyróżnik]
\label{aa:lem:kov2-trace-discriminant}
Niech $r,\bar r$ będą dwoma
różnymi sprzężonymi nad $K$
rozwiązaniami $Dr+r^2=q$,
gdzie $q\in K$.
Oznaczmy
$\phi=r+\bar r\in K$ oraz
$\Delta=(r-\bar r)^2\in K$.
Wówczas
\begin{equation}\label{aa:eq:kov2-discriminant}
 \Delta=4q-2D\phi-\phi^2,\qquad
 D\Delta=-2\phi\Delta,
\end{equation}
a ich wielomian minimalny ma postać
\begin{equation}\label{aa:eq:kov2-polynomial}
 T^2-\phi T+
 \frac{D\phi+\phi^2-2q}{2}.
\end{equation}
Ponadto $\Delta\ne0$ i
$\Delta$ nie jest kwadratem w $K$.
\end{lemat}
\begin{proof}
Ponieważ oba pierwiastki spełniają
Riccatiego,
\[
 D\phi=2q-r^2-\bar r^2
      =2q-\phi^2+2r\bar r,
\]
zatem
$r\bar r=(D\phi+\phi^2-2q)/2$.
Stąd
$\Delta=\phi^2-4r\bar r
=4q-2D\phi-\phi^2$
i otrzymujemy wielomian~\eqref{aa:eq:kov2-polynomial}.
Odejmując równania Riccatiego
dla obu pierwiastków,
otrzymujemy
$D(r-\bar r)=-(r+\bar r)(r-\bar r)$.
Różniczkowanie kwadratu daje
$D\Delta=-2\phi\Delta$.
Gdyby $\Delta$ była
kwadratem w $K$, równanie
\eqref{aa:eq:kov2-polynomial}
rozkładałoby się na czynniki
liniowe nad $K$, wbrew
$[K(r):K]=2$.
\end{proof}

\begin{lemat}[Wielomianowe równanie śladu]
\label{aa:lem:kov2-trace-equation}
Niech $q,\phi\in K$.
Jeżeli $\phi$ jest śladem dwóch
różnych rozwiązań Riccatiego,
to
\begin{equation}\label{aa:eq:kov2-trace}
 D^2\phi+3\phi D\phi+\phi^3
     -4q\phi-2Dq=0.
\end{equation}
Odwrotnie, jeśli $\phi\in K$
spełnia~\eqref{aa:eq:kov2-trace}, to
$\Delta=4q-2D\phi-\phi^2$
spełnia $D\Delta=-2\phi\Delta$.
Jeśli $\Delta\ne0$, dwa elementy
\begin{equation}\label{aa:eq:kov2-root}
 r_\pm=\frac{\phi\pm\sqrt{\Delta}}2
\end{equation}
w rozszerzeniu
$K(\sqrt{\Delta})$ spełniają
równanie Riccatiego.
Jeżeli $\Delta$ nie jest kwadratem
w $K$, oba rozwiązania są
ściśle kwadratowe i sprzężone.
\end{lemat}
\begin{proof}
Z pierwszej równości
\eqref{aa:eq:kov2-discriminant}
otrzymujemy po różniczkowaniu
\[
 D\Delta=4Dq-2D^2\phi-2\phi D\phi.
\]
Po dodaniu $2\phi\Delta$ prawa
strona wynosi
$-2(D^2\phi+3\phi D\phi+\phi^3
-4q\phi-2Dq)$.
Zatem warunek~\eqref{aa:eq:kov2-trace}
jest równoważny
$D\Delta=-2\phi\Delta$.
Załóżmy teraz $\Delta\ne0$
i wybierzmy $\delta^2=\Delta$.
W charakterystyce $0$
jedyna przedłużona derywacja
spełnia
$D\delta=(D\Delta)/(2\delta)
=-\phi\delta$.
Dla $r_\pm=(\phi\pm\delta)/2$ mamy
\[
 Dr_\pm+r_\pm^2
 =\frac{D\phi}{2}
   +\frac{\phi^2+\Delta}{4}=q;
\]
wyrazy zawierające $\delta$
znoszą się dzięki $D\delta=-\phi\delta$.
Dodatkowy warunek niekwadratowości
$\Delta$ jest dokładnie warunkiem
$[K(r_\pm):K]=2$.
\end{proof}

\begin{uwaga}[Weryfikacja negatywna]
\label{aa:uw:kov2-quadratic-versus-rational}
Równanie~\eqref{aa:eq:kov2-trace}
może mieć wymierne rozwiązanie
$\phi$ również wtedy, gdy
$\Delta$ jest kwadratem w $K$
albo gdy $\Delta=0$.
Wtedy
\eqref{aa:eq:kov2-root}
daje rozwiązania wymierne,
które należą już do
gałęzi~\ref{aa:tw:kov-rational-complete-algorithm}.
Dlatego sam test trzeciego
rzędu \emph{bez sprawdzenia
wyróżnika} nie dowodzi,
że otrzymaliśmy nowy
przypadek kwadratowy.
\end{uwaga}

\section{Symetryczny kwadrat równania liniowego}

\begin{lemat}[Równanie trzeciego rzędu]
\label{aa:lem:kov2-symmetric-square}
Niech $q$ należy do
ciała różniczkowego charakterystyki
$0$, a $z_1,z_2$ spełniają
$D^2z_j=qz_j$ w jego
rozszerzeniu różniczkowym.
Wówczas ich iloczyn
$v=z_1z_2$ spełnia
\begin{equation}\label{aa:eq:kov2-third-order}
 D^3v-4qDv-2(Dq)v=0.
\end{equation}
Jeżeli $v\ne0$, to
$\phi=Dv/v$ spełnia
\eqref{aa:eq:kov2-trace}.
\end{lemat}
\begin{proof}
Różniczkując iloczyn,
dostajemy
$D^2v=2qv+2(Dz_1)(Dz_2)$.
Po kolejnym różniczkowaniu:
\[
 D^3v=2(Dq)v+2qDv
 +2q\bigl(z_1Dz_2+(Dz_1)z_2\bigr)
 =2(Dq)v+4qDv.
\]
Ponadto z $Dv=\phi v$ wynika
$D^3v=(D^2\phi+3\phi D\phi
+\phi^3)v$.
Po podstawieniu do poprzedniej
równości i podzieleniu
przez $v$ otrzymujemy
\eqref{aa:eq:kov2-trace}.
\end{proof}

\begin{lemat}[Liniowe równanie dla wielomianu]
\label{aa:lem:kov2-polynomial-equation}
Niech $q,\theta\in\mathbb C(x)$,
a $P\in\mathbb C[x]$ będzie
niezerowy. Połóżmy
$\phi=\theta+P'/P$.
Wówczas $\phi$ spełnia
\eqref{aa:eq:kov2-trace}
wtedy i tylko wtedy, gdy
\begin{equation}\label{aa:eq:kov2-P-equation}
\begin{split}
 P''' +3\theta P''
  +(3\theta^2+3\theta'-4q)P'
  +\bigl(\theta''+3\theta\theta'
         +\theta^3-4q\theta-2q'\bigr)P=0.
\end{split}
\end{equation}
Po pomnożeniu przez
wspólny mianownik jest to
\emph{liniowy} układ równań
na współczynniki $P$
ustalonego stopnia.
\end{lemat}
\begin{proof}
Niech $v$ spełnia $Dv=\theta v$,
$v\ne0$, i rozważmy $w=Pv$.
Wtedy $Dw/v=P'+\theta P$,
a iterując otrzymujemy
\[
 D^3w/v=P'''+3\theta P''
 +(3\theta^2+3\theta')P'
 +(\theta''+3\theta\theta'+\theta^3)P.
\]
Po podstawieniu w równanie
\eqref{aa:eq:kov2-third-order}
widzimy, że $w$ jest jego
rozwiązaniem wtedy i tylko wtedy,
gdy zachodzi
\eqref{aa:eq:kov2-P-equation}.
Ponieważ
$Dw/w=\theta+P'/P=\phi$,
równoważność z
\eqref{aa:eq:kov2-trace}
wynika z poprzedniego lematu
i tożsamości na pochodne
logarytmiczne. Wszystkie
współczynniki równania
\eqref{aa:eq:kov2-P-equation}
są ustalonymi funkcjami
wymiernymi, więc porównanie
współczynników po usunięciu
mianowników daje liniowy
układ algebraiczny.
\end{proof}

\section{Klasyczne lokalne listy Kovacica}

\begin{definicja}[Rząd w nieskończoności]
\label{aa:def:kov2-infinity-order}
Niech $q=p/s\in\mathbb C(x)$,
gdzie $p,s\in\mathbb C[x]$
są względnie pierwsze,
$s\ne0$, $p\ne0$.
\emph{Rzędem $q$ w nieskończoności}
nazywamy liczbę
$\operatorname{ord}_{\infty}q
=\deg s-\deg p$.
Jeżeli $\operatorname{ord}_{\infty}q=2$,
przez $b_\infty$ oznaczamy
współczynnik przy $x^{-2}$
w rozwinięciu Laurenta $q$
w nieskończoności.
Dla $q=0$ stosujemy
najpierw gałąź wymierną
z rozdziału~\ref{aa:ch:kovacic-rational-potential}.
\end{definicja}

\begin{twierdzenie}[Listy lokalne przypadku drugiego -- wynik klasyczny]
\label{aa:tw:kov2-classical-candidates}
Niech $q\in\mathbb C(x)$ będzie
niezerowe, a $\Gamma$ skończonym
zbiorem wszystkich jego biegunów
w $\mathbb C$.
Dla każdego $a\in\Gamma$
definiujemy skończony zbiór
$E_a\subseteq\mathbb Z$:
\begin{enumerate}
\item jeśli biegun ma rząd $1$,
to $E_a=\{4\}$;
\item jeśli biegun ma rząd $2$,
a $b_a$ jest współczynnikiem
przy $(x-a)^{-2}$, to
\[
 E_a=\{2,\;2+2\sqrt{1+4b_a},
                 \;2-2\sqrt{1+4b_a}\}
                 \cap\mathbb Z;
\]
\item jeśli biegun ma rząd $m>2$,
to $E_a=\{m\}$.
\end{enumerate}
Również w nieskończoności
definiujemy zbiór:
\[
 E_\infty=
 \begin{cases}
 \{0,2,4\},
   &\operatorname{ord}_\infty q>2,\\
 \{2,\;2+2\sqrt{1+4b_\infty},
              \;2-2\sqrt{1+4b_\infty}\}
              \cap\mathbb Z,
   &\operatorname{ord}_\infty q=2,\\
 \{\operatorname{ord}_\infty q\},
   &\operatorname{ord}_\infty q<2.
 \end{cases}
\]
Jeśli równanie $r'+r^2=q$
ma ściśle kwadratowe
rozwiązanie nad $\mathbb C(x)$,
to istnieją wybory
$e_a\in E_a$ oraz
$e_\infty\in E_\infty$
takie, że
\begin{equation}\label{aa:eq:kov2-degree}
 n=\frac12\left(e_\infty-
                      \sum_{a\in\Gamma}e_a\right)
 \in\mathbb Z_{\ge0}
\end{equation}
oraz dla
\begin{equation}\label{aa:eq:kov2-theta}
 \theta=\frac12\sum_{a\in\Gamma}
                     \frac{e_a}{x-a}
\end{equation}
istnieje unormowany
$P\in\mathbb C[x]$
stopnia $n$ spełniający
\eqref{aa:eq:kov2-P-equation}.
\end{twierdzenie}

\begin{uwaga}[Status dowodu list lokalnych]
\label{aa:uw:kov2-local-cited}
To twierdzenie cytujemy
w standardowej postaci
drugiego przypadku algorytmu
Kovacica: J.~Kovacic,
\emph{An algorithm for solving
second order linear homogeneous
differential equations},
\emph{Journal of Symbolic Computation}
2 (1986), 3--43.
Dla biegunów drugiego rzędu
liczby $2\pm2\sqrt{1+4b_a}$
wynikają z lokalnych
wykładników równania
$z''=qz$, ale pełny
dowód wyczerpującej listy
dla wszystkich rzędów,
w tym rozgałęzienia
w punktach osobliwych,
wymaga analizy lokalnej
użytej przez Kovacica.
Właśnie ta konieczność
i wyczerpujący charakter
list $E_a,E_\infty$
\emph{nie są tutaj
samodzielnie dowodzone}.
Resztę algorytmu
wyprowadzamy i sprawdzamy
wprost z
Lematów~\ref{aa:lem:kov2-trace-equation}
i~\ref{aa:lem:kov2-polynomial-equation}.
\end{uwaga}

\section{Skończona procedura drugiej gałęzi}

\begin{twierdzenie}[Druga gałąź Kovacica, z cytowanym twierdzeniem lokalnym]
\label{aa:tw:kov2-algorithm}
Niech $q\in\mathbb C(x)$,
a jego współczynniki
będą efektywnie zadane
w sensie
Uwagi~\ref{aa:uw:kov-rational-effective}.
Przyjmując klasyczne
Twierdzenie~\ref{aa:tw:kov2-classical-candidates},
istnieje skończony algorytm
rozstrzygający, czy
$r'+r^2=q$ ma
rozwiązanie ściśle kwadratowe
nad $\mathbb C(x)$,
i konstruujący takie
rozwiązanie, jeśli istnieje.
\end{twierdzenie}
\begin{proof}
Dla $q=0$ dowolne rozwiązanie
wymierne zapewnia gałąź
\ref{aa:tw:kov-rational-complete-algorithm};
ewentualne przypadki
ściśle kwadratowe
wykluczamy bezpośrednio:
jeżeli $r'+r^2=0$ i $r\ne0$,
to $D(1/r)=1$,
stąd $1/r=x+c$, gdzie
$c$ jest stałą w algebraicznym
rozszerzeniu $\mathbb C(x)$.
Ponieważ $\mathbb C$ jest
algebraicznie domknięte,
$c\in\mathbb C$, a więc
$r=1/(x+c)\in\mathbb C(x)$.
Dla $r=0$ również
nie ma stopnia drugiego.

Niech dalej $q\ne0$.
Wyznaczamy bieguny
$q$, ich rzędy oraz
$\operatorname{ord}_\infty q$.
Tworzymy skończone
listy $E_a,E_\infty$
z cytowanego twierdzenia.
Rozważamy wszystkie
skończenie liczne kombinacje
$(e_\infty,(e_a)_{a\in\Gamma})$;
odrzucamy te, dla których
$n$ określone w~\eqref{aa:eq:kov2-degree}
nie należy do
$\mathbb Z_{\ge0}$.

Dla każdej pozostałej
kombinacji obliczamy $\theta$
z~\eqref{aa:eq:kov2-theta}
i podstawiamy
$P=x^n+\sum_{j<n}p_jx^j$
do równania
\eqref{aa:eq:kov2-P-equation}.
Po pomnożeniu przez wspólny
mianownik otrzymujemy
skończony liniowy układ
nad efektywnym ciałem
współczynników. Jeśli nie
ma rozwiązań, odrzucamy
daną kombinację.
Jeśli ma, jego zbiór
rozwiązań stanowi
afiniczną przestrzeń $V$.

Dla każdego $P\in V$
obliczamy
$\phi=\theta+P'/P$
oraz
$\Delta=4q-2\phi'-\phi^2$.
Jeżeli $\Delta\ne0$
i nie jest kwadratem
w $\mathbb C(x)$,
Lemat~\ref{aa:lem:kov2-trace-equation}
daje ściśle kwadratowe
rozwiązania
$r_\pm=(\phi\pm\sqrt{\Delta})/2$.
Test kwadratowości funkcji
wymiernej jest efektywny:
w jej zredukowanym
rozkładzie na
czynniki wykładniki
wszystkich czynników
muszą być parzyste;
ponieważ $\mathbb C$
jest algebraicznie domknięte,
stały współczynnik
nie stanowi przeszkody.

Pozostaje uzasadnić,
że istnienie odpowiedniego
$P\in V$ można
rozstrzygnąć bez
nieskończonego przeszukiwania.
Po podstawieniu
$P=\sum p_jx^j$
i usunięciu mianowników
warunek $\Delta=0$
albo
$\Delta$ jest kwadratem
w $\mathbb C(x)$
opisuje algebraiczne
warunki na parametry
$P$ (można użyć
rozkładu wielomianów,
NWD i kwantyfikatorów
nad algebraicznie
domkniętym ciałem).
Dla ustalonego stopnia $P$
jest to skończony problem
eliminacji algebraicznej:
rozstrzyga się, czy
afininiczny zbiór $V$
zawiera punkt, dla
którego $\Delta$ jest
niezerowa i niewymiernie
pierwiastkowalna.
Równoważnie można
skorzystać z pełnego
rozstrzygania zdań
w teorii ciał algebraicznie
domkniętych charakterystyki
$0$; w warunkach efektywnej
prezentacji współczynników
jest to procedura
kończąca się.

\emph{Dostateczność:}
każde znalezione $P$,
po obliczeniu wyróżnika
i wybraniu pierwiastka,
daje dokładnie rozwiązanie
Riccatiego i lokalne
rozwiązanie $z''=qz$
postaci
$z=\exp(\int r\,\mathrm dx)$,
na mocy
Lematu~\ref{aa:lem:kov2-trace-equation}.
\emph{Konieczność:}
jeśli istnieje rozwiązanie
ściśle kwadratowe,
to klasyczne
Twierdzenie~\ref{aa:tw:kov2-classical-candidates}
zapewnia jedną z
badanych kombinacji
i pewien unormowany
wielomian $P$
spełniający układ.
Jego ślad $\phi$
ma wyróżnik niekwadratowy,
a zatem jest dopuszczony
przez końcowy test.
To dowodzi skończoności
oraz obu kierunków
\emph{przy wyraźnie
cytowanym twierdzeniu lokalnym}.
\end{proof}

\begin{uwaga}[O kosztach i implementacji]
\label{aa:uw:kov2-computation}
Dla danych z prostymi
parametrami w większości
przykładów liniowy układ
wyznacza jedyne $P$.
Dopuszczenie ogólnych
parametrycznych rozwiązań
i rozstrzyganie
kwadratowości wyróżnika
wymaga dodatkowej kontroli,
nie zaś arbitralnego
wyboru pierwszego
rozwiązania układu.
Użycie eliminacji
algebraicznej zapewnia
teoretyczną skończoność,
ale nie jest twierdzeniem
o praktycznej wydajności.
W szczególności
\emph{testy symboliczne
nie zastępują tego dowodu}.
\end{uwaga}

\section{Przykłady z pełnym rachunkiem}

\begin{przyklad}[Kwadratowy pierwiastek i wielomian stopnia zero]
\label{aa:prz:kov2-zero-degree}
Rozważmy
\[
 q=x+\frac5{16x^2}.
\]
W punkcie $0$ występuje
biegun drugiego rzędu
o $b_0=5/16$.
Stąd
$\sqrt{1+4b_0}=3/2$
i $E_0=\{-1,2,5\}$.
W nieskończoności
$\operatorname{ord}_\infty q=-1$,
więc $E_\infty=\{-1\}$.
Wybierając $e_0=-1$
i $e_\infty=-1$,
otrzymujemy $n=0$ oraz
\[
 \theta=-\frac1{2x},\qquad P=1.
\]
Bezpośrednie podstawienie
do równania
\eqref{aa:eq:kov2-P-equation}
daje zero. Zatem
$\phi=-1/(2x)$ oraz
\[
 \Delta=4q-2\phi'-\phi^2=4x.
\]
Nie jest to kwadrat w
$\mathbb C(x)$.
Otrzymujemy dokładnie
\[
 r_\pm=-\frac1{4x}\pm\sqrt{x},
 \qquad r_\pm'+r_\pm^2=q.
\]
Odpowiadające
lokalne rozwiązania
to
$x^{-1/4}\exp(\pm\tfrac23 x^{3/2})$.
\end{przyklad}

\begin{przyklad}[Niezerowy stopień wielomianu pomocniczego]
\label{aa:prz:kov2-positive-degree}
Niech
\[
 q=\frac{4x+5}{16x^2}
   =\frac1{4x}+\frac5{16x^2}.
\]
Przy $x=0$ znów
$E_0=\{-1,2,5\}$,
lecz teraz
$\operatorname{ord}_\infty q=1$
i $E_\infty=\{1\}$.
Dla $e_0=-1$,
$e_\infty=1$
wzór~\eqref{aa:eq:kov2-degree}
daje $n=1$, a
$\theta=-1/(2x)$.
Podstawiamy
$P=x+p_0$ do równania
\eqref{aa:eq:kov2-P-equation}
i porównujemy współczynniki.
Otrzymujemy $p_0=-1$,
czyli $P=x-1$.
W szczególności
\[
 \phi=-\frac1{2x}+\frac1{x-1}
      =\frac{x+1}{2x(x-1)},
 \qquad
 \Delta=\frac{x}{(x-1)^2}.
\]
Wyróżnik nie jest
kwadratem w $\mathbb C(x)$,
ponieważ $x$ ma
w $0$ nieparzystą
krotność zera.
Pierwiastki Riccatiego wynoszą
\[
 r_\pm=-\frac1{4x}
       +\frac1{2(x-1)}
       \pm\frac{\sqrt{x}}{2(x-1)}.
\]
Podstawienie
$r_\pm'+r_\pm^2=q$
jest bezpośrednim
sprawdzeniem wyniku.
Przykład pokazuje,
że pominięcie $P$
pozwalałoby znaleźć
tylko przypadek $n=0$.
\end{przyklad}

\begin{przyklad}[Inna konfiguracja lokalna]
\label{aa:prz:kov2-infinity-one}
Dla
$q=1/x-3/(16x^2)$
mamy
$b_0=-3/16$,
$\sqrt{1+4b_0}=1/2$
i $E_0=\{1,2,3\}$.
W nieskończoności
$\operatorname{ord}_\infty q=1$,
więc $e_\infty=1$.
Wybór $e_0=1$
daje $n=0$,
$\theta=1/(2x)$
i $P=1$.
Następnie
\[
 \Delta=4q-2\theta'-\theta^2=\frac4x,
 \qquad
 r_\pm=\frac1{4x}\pm\frac1{\sqrt{x}}.
\]
Każdy pierwiastek
spełnia Riccatiego,
co można sprawdzić
przez różniczkowanie.
Zatem istnieją lokalne
rozwiązania
$z=x^{1/4}\exp(\pm2\sqrt{x})$.
\end{przyklad}

\begin{przyklad}[Wynik negatywny dla równania Airy'ego]
\label{aa:prz:kov2-airy}
Dla $q=x$ nie ma
biegunów skończonych.
Ponieważ
$\operatorname{ord}_\infty q=-1$,
mamy $E_\infty=\{-1\}$,
a więc
$n=-1/2\notin\mathbb Z_{\ge0}$.
Lista drugiej gałęzi
jest pusta.
W rozdziale~\ref{aa:ch:picard-vessiot}
udowodniono silniejszy
wynik: równanie Airy'ego
nie ma \emph{żadnych}
niezerowych rozwiązań
liouvillowskich.
Sam obecny test
dotyczy tylko drugiej
gałęzi algorytmu.
\end{przyklad}

\section{Związek z grupą Galois}

Niech $r_\pm$
będą ściśle kwadratowymi,
sprzężonymi rozwiązaniami
Riccatiego.
Odpowiadające im
rozwiązania $z_\pm$
równania $z''=qz$
wyznaczają dwie różne
proste w dwuwymiarowej
przestrzeni rozwiązań.
Automorfizm różniczkowy
nad $\mathbb C(x)$
przekształca zbiór
$\{r_+,r_-\}$
w siebie, zatem
przekształca także
parę tych prostych
w siebie. Grupa Galois
zachowuje więc
\emph{nieuporządkowaną}
parę prostych,
co odpowiada
podgrupie normalizatora
torusa diagonalnego
w $\mathrm{SL}_2(\mathbb C)$.
Jeśli odrzuciliśmy
wcześniej przypadek
wymierny z AA-48,
nie istnieje
niezmiennicza pojedyncza
prosta zdefiniowana
nad $\mathbb C(x)$.
Powyższe rozumowanie
jest ilustracją
imprymitywnego
przypadku teorii Galois;
pełna klasyfikacja
grup i odpowiedniość
Galois zostały
sformułowane jako
wyniki klasyczne
w rozdziale~\ref{aa:ch:picard-vessiot}.

Trzeci przypadek algorytmu
Kovacica, związany
z grupami skończonymi
i stopniami algebraicznymi
$4$, $6$ i $12$,
nie należy do zakresu
AA-49. Jego
osobne opracowanie wymaga
nowego rozdziału
i analogicznych audytów.
